\documentclass[10pt,a4paper]{amsart}
\usepackage[margin=19mm]{geometry}
\usepackage{amsmath,amssymb,mathtools}
\usepackage[scr=rsfs,cal=euler]{mathalfa}
\usepackage{xfrac}
\usepackage[unicode,hidelinks,bookmarks=false]{hyperref}
\newcommand{\ie}{i.e.\ }
\newcommand{\cf}{cf.\ }
\newcommand{\resp}{respectively\ }
\newcommand{\Subsection}{\subsection}
\providecommand{\nobreakdash}{\nobreak\hbox{-}}
\setlength{\emergencystretch}{2em}
\multlinegap0pt
\usepackage{mathtools}
\usepackage{amssymb}
\usepackage{xcolor}
\newtheorem{lemma}{Lemma}
\title{A Bijection Between Stacked Directed Polyominoes \\ and Motzkin Paths with Catastrophes}
\author{Florian Schager, Michael Wallner}
\date{Source: arXiv:2406.16417v2 (CC BY 4.0)}
\begin{document}
\maketitle

\noindent\emph{Source and licence.} A Bijection Between Stacked Directed Polyominoes \\ and Motzkin Paths with Catastrophes. Version arXiv:2406.16417v2 (CC BY 4.0), distributed by arXiv under the Creative Commons Attribution 4.0 International licence, http://creativecommons.org/licenses/by/4.0/, as recorded in the arXiv licence metadata for this exact version and shown on the abstract page https://arxiv.org/abs/2406.16417v2. Full licence notice: \texttt{licenses/arxiv-2406.16417-CC-BY-4.0.txt}.\par\smallskip

\noindent\textit{Editorial scope.} Counted unit: the article's lemma lem:no-SE (source0.tex lines 603--607). The article's complete original proof of it is delivered with it (source0.tex lines 609--625, one \texttt{proof} environment of 17 lines). No mathematics is written, repaired or re-derived: the statement and the proof are the article's own lines, in the article's own notation, and no retained line is substituted or re-flowed.  No further source-local context is needed: every cross-reference of every retained line resolves inside the two delivered spans.  Nothing was omitted from any retained span, so the package carries no omission record.  No retained line contains a citation, so the wrapper carries no bibliography and no bibliography entry.  No retained line contains a figure environment, an image-inclusion command or drawing code, so no image is delivered and the package declares no asset dependency.  Editorial formatting only, in addition: an AMS \texttt{amsart} preamble that restates the article's own declarations and macros used by the retained lines, namely the environments \texttt{lemma} on separate counters, together with the standard packages those lines need.  The retained environments print in this document as Lemma 1 in order of appearance, whereas the article prints the same items with the numbering of its own full document.  The complete native source of the article as distributed in the arXiv e-print for this exact version is bundled unchanged as \texttt{sources/161263/source0.tex}.
\begin{lemma}
  \label{lem:no-SE}
  Let $X_n$ denote the random variable counting the number of \textbf{SE} steps in a Motzkin excursion of length $n$ chosen uniformly at random.
  Then, $\mathbb{E}[X_n] \xrightarrow[]{{n \to \infty}} \frac{n}{7}$.
\end{lemma}

\begin{proof}
  We mark the \textbf{SE} steps in a Motzkin path by introducing the bivariate jump polynomial $P(u,x) := \frac{x}{u} + 1 + u$. 
  This leads to the bivariate generating function
  \[
  E_\mathcal{M}(z,x) = \frac{u_1(z,x) \left(-1+u_1(z,x) \right)}{z \left(-1+\left(x -u_1(z,x) +3\right) z \right) x},
  \]
  of Motzkin excursion with catastrophes, where $u_1(z,x)$ denotes the small branch of the kernel equation $1 - z \cdot P(u,x) = 0$. The variable $x$ counts the number of \textbf{SE} steps.
   Then, 
  \[
    \mathbb{P}[X_n = k] = \frac{[z^n x^k] E_\mathcal{M}(z,x)}{[z^n] E_\mathcal{M}(z,1)} \quad \text{and} \quad \mathbb{E}[X_n] = [z^n] \left(\frac{\frac{\partial}{\partial x} E_\mathcal{M}(z,1)}{ E_\mathcal{M}(z,1)}\right).
  \]
  With the help of a computer algebra system one obtains that
  \[
    [z^n] \frac{\partial}{\partial x} E_\mathcal{M}(z,1) = \frac{3n}{56} \cdot \left( \frac{7}{2}\right)^n \cdot \left(1 + \mathcal{O}\left(\frac{1}{n}\right)\right)  \quad \text{and} \quad [z^n] E_\mathcal{M}(z,1) = \frac{3}{8} \cdot \left( \frac{7}{2}\right)^n \cdot \left(1 + \mathcal{O}\left(\frac{1}{n}\right)\right)
  \]
  and thus $\mathbb{E}[X_n] \xrightarrow[]{{n \to \infty}} \frac{n}{7}$.
\end{proof}

\end{document}
